% ═════════════════════════════════════════════════════════════════════════════
\section{Early Developments}
\label{sec:early}
% ═════════════════════════════════════════════════════════════════════════════

The questions studied in this chapter were first formulated in a body of work
from the late 1980s and 1990s that established both the model and the basic
techniques of parameterized verification. This section reviews that work and
some of its later developments.

\paragraph{Correspondence between instances.}
While the undecidability results of Section~\ref{sec:prelim} showed early on
that a general algorithm cannot exist~\cite{AK86,Suz88}, Browne, Clarke, and
Grumberg~\cite{BCG89} proposed a way of transferring
verification results from small instances to large ones. They introduced an
indexed version of the temporal logic CTL$^*$ without the next-time operator,
together with a notion of \emph{correspondence} between systems of different
sizes that preserves the formulas of this logic. Once a correspondence between
a small instance and all larger instances has been established, model checking
the small instance suffices. Establishing the correspondence requires, in
general, an argument by the user, but the approach anticipates the cutoff
theorems discussed in Section~\ref{sec:other}, in which the transfer from small to large
instances is guaranteed by the structure of the system.

\paragraph{Systems with many processes.}
A foundational formulation considers a system composed of a single control
process together with an unbounded collection of user processes sharing a
common definition, communicating through rendezvous synchronization in the
style of CCS~\cite{GerSis92}. The central contribution was to show that
reasoning about the entire infinite family can be reduced to a finite analysis:
decision procedures were given for checking whether every process satisfies a
given indexed temporal specification. When all processes are identical, the
procedure runs in time polynomial in the size of the process template; in the
presence of a control process, it runs in doubly exponential time. The
reachability questions arising in the analysis were connected to the theory of
Petri nets and vector addition systems, through a \emph{counting abstraction}:
since the user processes are anonymous, it
suffices to record how many of them are in each local state. The question of
whether some user processes can reach a given combination of local states
 becomes the
\emph{coverability} problem for Petri nets, which is decidable and, by the
results of Lipton~\cite{Lip76} and Rackoff~\cite{Rac78},
\textsc{ExpSpace}-complete~\cite{Esp98}.
This early demarcation of the boundary between decidable and undecidable cases,
and of the algebraic structures underlying it, shaped much of the work that
followed.

\paragraph{Induction over the number of processes.}
A second early line of work used induction on the number of processes. Kurshan
and McMillan~\cite{KM89}, and independently Wolper and
Lovinfosse~\cite{WL89}, observed that if a process~$I$ is found such that $I$
\emph{abstracts} a single copy of the template, in a sense that preserves the
properties of interest, and the composition of~$I$ with one more copy is again
abstracted by~$I$, then $I$ abstracts the composition of any number of copies.
The
process~$I$ is then a \emph{network invariant}, and verifying the composition
of~$I$ with a fixed environment is sufficient for all sizes. We return to this
idea in Section~\ref{sec:other}. Clarke, Grumberg, and
Jha~\cite{CGJ97} generalized the setting to families of networks generated by
\emph{network grammars}, using abstraction to make the search for invariants
effective.

A complementary line of work targeted linear networks of arbitrary size and
introduced automata as inductive invariants~\cite{SisGyu99}. A correctness
property is itself specified by an automaton, from which an invariant
automaton is synthesized; an induction on the number of processes then
establishes that the property is preserved as the network grows. The method
was shown to be complete for safety properties. It is closely related to both
the network-invariant approach of Section~\ref{sec:other} and the regular
model checking framework of Section~\ref{sec:rmc}, in which sets of global
configurations of unbounded width are represented symbolically by finite
automata.

\paragraph{Symmetry.}
A further theme running through this early work is the exploitation of the
symmetry inherent in a population of identical processes. Because the identity
of an individual process is immaterial to the global behavior of the system,
verification can be carried out over equivalence classes of configurations
rather than over configurations themselves~\cite{EmeSis96,CEFJ96,ID96}. Two
configurations are equivalent if one is obtained from the other by permuting
the process indices, and the quotient of the state space under this
equivalence can be exponentially smaller than the original. Symmetry
reduction does not by itself eliminate the need to reason about an arbitrary number of processes, but it explains why configurations of symmetric
systems can be reduced to simple multisets of local process states, and it is the basis of the counting
abstraction above. Combining symmetry with fairness leads to on-the-fly model
checking procedures that keep the explored state space
manageable~\cite{GyuSis97}.

\paragraph{Broadcast protocols and cache coherence.}
Esparza, Finkel, and Mayr~\cite{EFM99} studied \emph{broadcast protocols}, in
which a process can send a message that is received by all other processes in
the same step. Broadcast protocols are an extension of Petri nets, in which a
transition can move all tokens from one place to another. Safety properties
expressed by upward-closed sets remain decidable for broadcast protocols, and
can be checked by backward reachability analysis (Section~\ref{sec:wqo}), while
liveness properties become undecidable. Delzanno~\cite{Del00} applied this
framework to cache coherence protocols. Using the counting abstraction and
symbolic representations based on linear arithmetic constraints, he verified
several coherence protocols automatically for an arbitrary number of caches.
This was one of the first demonstrations that parameterized verification could
handle protocols of practical relevance.

\paragraph{Later developments.}
These ideas have been extended in many directions. Parameterized systems with
pairwise rendezvous have been analyzed on a range of
topologies~\cite{AKRSV18}. In \emph{asynchronous shared-memory systems}, an
unbounded number of anonymous processes communicate through shared registers
without locks~\cite{EGM16}. In \emph{population protocols}, anonymous agents
interact in pairs to compute a predicate of their initial
states~\cite{AADFP06}, and verifying that a protocol computes the intended
predicate for every population size is a parameterized problem that can be
reduced to problems on Petri nets~\cite{EGLM17}. In \emph{ad hoc networks},
processes communicate by broadcast to their neighbors in an unknown graph; the
problem is undecidable in general, and becomes decidable under restrictions on
the possible topologies~\cite{DSZ10}. A recurring observation in this line of
work is that, when processes are anonymous and any behavior of one process can
be replicated by others, the complexity of parameterized verification is often
much lower than that of Petri net coverability~\cite{Esp14}.

Taken together, these strands form the conceptual foundation on which the
remainder of this chapter builds: finite reductions for synchronizing
populations, induction over the number of processes, automata-based
invariants for linear networks, symmetry reduction for replicated components,
and well-structured models of broadcast.
